\subsection{Профессор Хаос}

\subsubsection{Описание решения}

В задаче для числа дней эксперимента используется ограничение \( 1\leq k\leq 10^{18} \), что намекает на математическую природу решения. Линейный алгоритм для моделирование задачи будет бессилен при больших \( k \).

Опуская детали, на каждом дне эксперимента к количеству бактерий применяется линейная функция \( f(x)=bx-c \). Заметим, что если \( c=a\cdot(b-1) \), то \( f(a)=a \) в любой из дней эксперимента \textit{(\( a \) --- неподвижная точка \( f \))}.

Если после применения \( f \) число бактерий \textit{увеличилось}, то на следующем шаге оно также будет увеличиваться в силу возрастания \( f \). Но по условию задачи число бактерий ограничено \( d \), поэтому при достижении этой отметки число бактерий останется неизменным в последующих днях эксперимента.

Если после применения \( f \) число бактерий \textit{уменьшилось}, то на следующем шаге оно также будет уменьшаться в силу возрастания \( f \). Значит, оно в какой-то из дней достигнет нуля.

\subsubsection{Асимптотическая сложность}

\textit{Временная сложность решения} --- \( \mathcal{O}(d) \). Если число бактерий не изменяется, ответ однозначен. Иначе будет выполнено не более \( d \) операций: функция целая, и она в худшем случае может расти или убывать на единицу до тех пор, пока не достигнет границ отрезка \( [0,d] \).

\textit{Пространственная сложность решения} --- \( \mathcal{O}(1) \). В решении используются только переменные для хранения промежуточных данных и ввода/вывода.

\subsubsection{Листинг кода}

\begin{lstlisting}
#include <iostream>

int main() {
  int a, b, c, d;
  long long k;
  std::cin >> a >> b >> c >> d >> k;
  if (c == a * (b - 1)) {
    std::cout << a;
    return 0;
  }
  long x = a;
  for (int i = 0; i < (d < k ? d : k); i++) {
    x *= b;
    if (x < c) {
      x = 0;
      break;
    }
    x -= c;
    if (x >= d) {
      x = d;
      break;
    }
  }
  std::cout << x;
  return 0;
}
\end{lstlisting}